\documentclass[11pt]{article}
\usepackage[margin=1in]{geometry}
\usepackage{amsmath,amssymb,amsfonts}
\usepackage{bm}
\usepackage{enumitem}
\usepackage{hyperref}
\hypersetup{colorlinks=true,linkcolor=blue,urlcolor=blue}

\newcommand{\Lam}{\Lambda}
\newcommand{\X}{\mathcal{X}}
\newcommand{\Hist}{\mathcal{H}}
\newcommand{\softplus}{\operatorname{softplus}}

\title{\textbf{LGM: Linear Ground-rate \texttimes\ deep soft-max Marks}\\[2pt]
\large A rate-pinned set-valued marked temporal point process for limit-order-book event streams}
\author{}
\date{}

\begin{document}
\maketitle
\vspace{-3em}

\section{Event stream and factorization}
We model the order-book event stream as a sequence of \emph{set-valued} continuous-time events
\[
\Hist_T=\{(t_i,x_i)\}_{i=1}^{N_T},\qquad
x_i\in\X:=\mathcal{P}(\{1,\dots,K\}),
\]
where $t_i$ is an arrival time and $x_i$ is the subset of the $K$ event types that occur
simultaneously at $t_i$. Following the set-valued MTPP formulation
(Chang, Boyd, Smyth, 2024) we factor the set-specific intensity into a scalar
arrival rate and a normalized mark distribution:
\[
\boxed{\;\lambda_x^*(t)=\Lam(t)\,p^*(x\mid t),\qquad x\in\X.\;}
\]
Because $\sum_{x\in\X}p^*(x\mid t)=1$, the mark head is \emph{rate-neutral}:
\[
\sum_{x\in\X}\lambda_x^*(t)=\Lam(t)\sum_{x\in\X}p^*(x\mid t)=\Lam(t).
\]
LGM specifies $\Lam(t)$ as a \textbf{linear multi-timescale Hawkes ground rate} (\S\ref{sec:ground})
and $p^*$ as a \textbf{deep soft-max / Bernoulli mark head} (\S\ref{sec:mark}); the two are coupled
only through the shared history embedding.

\section{Linear multi-timescale Hawkes ground rate}\label{sec:ground}
The ground intensity counts \emph{set-event arrivals} (one increment per timestamp $t_i$,
regardless of $|x_i|$) and is a linear Hawkes process with an exponential-mixture kernel:
\[
\Lam(t)=\mu_0+\int_0^t\varphi(t-u)\,dN(u),\qquad
\varphi(\tau)=\sum_{m=1}^{M}a_m e^{-\delta_m\tau},\quad a_m\ge 0.
\]
With the decayed traces $s_m(t)=\sum_{t_i<t}e^{-\delta_m(t-t_i)}$ this is the closed form
\[
\Lam(t)=\mu_0+\sum_{m=1}^{M}a_m\,s_m(t),
\]
which updates in $O(M)$ per event. Non-negativity is enforced by
$a_m=\softplus(\tilde a_m)$, and the $M$ decays $\delta_m$ span fast-to-slow timescales.

\section{Stationary-mean calibration (the rate pin)}\label{sec:pin}
Let $\bar\Lam=\mathbb{E}[\Lam(t)]$. Under stationarity $\mathbb{E}[dN(u)]=\bar\Lam\,du$, so taking
expectations of the intensity equation gives the fixed point
\[
\bar\Lam=\mu_0+\bar\Lam\,n,\qquad
n:=\int_0^\infty\varphi(\tau)\,d\tau=\sum_{m=1}^{M}\frac{a_m}{\delta_m}\ \ (\text{branching ratio}),
\]
hence $\bar\Lam=\mu_0/(1-n)$ for $n<1$. \textbf{Pinning} the base rate to
\[
\boxed{\;\mu_0=R_{\mathrm{target}}\,(1-n)\;}
\qquad\Longrightarrow\qquad
\bar\Lam=\frac{R_{\mathrm{target}}(1-n)}{1-n}=R_{\mathrm{target}}
\]
fixes the stationary mean arrival rate \emph{exactly} to $R_{\mathrm{target}}$, independent of the
learned clustering shape. Each forward pass computes
$n=\sum_m \softplus(\tilde a_m)/\delta_m$ (projected/clipped below a cap $n<1$) and sets
$\mu_0$ deterministically. Consequently the optimizer can learn temporal clustering through
$a_m/\delta_m$ but \emph{cannot} inflate the long-run event rate, and the subcritical
certificate $n<1$ is available in closed form (\texttt{closed\_form\_rho}).

\paragraph{Optional volatility feedback (heavy tails).}
A QHawkes-style mean-zero quadratic term on a learned signed order-flow state $R(t)$ can be added,
\[
\Lam(t)=\mu_0+\sum_m a_m s_m(t)+b\bigl(R(t)^2-c_0\bigr),\qquad c_0\approx\mathbb{E}[R^2],
\]
so the rate spikes during directional runs (producing fat tails / volatility clustering) while the
mean correction $c_0$ preserves the exact pin $\bar\Lam=R_{\mathrm{target}}$.

\section{Mark head: single-item vs.\ set-valued}\label{sec:mark}
Both heads consume the same per-type logits $z_k=z_k(h_\theta(t))$, produced by a
\emph{per-type parallel} continuous-time recurrent network (each type carries its own decaying
latent with a weight-shared read-out), and differ only in normalization and likelihood.

\paragraph{Single-item (categorical) head.} Support restricted to singletons $\{k\}$:
\[
p^*(\{k\}\mid t)=\frac{e^{z_k}}{\sum_{j=1}^{K}e^{z_j}},\qquad
\sum_{k=1}^{K}p^*(\{k\}\mid t)=1,
\]
trained by categorical cross-entropy. Each arrival is one event type.

\paragraph{Set-valued (Dynamic Bernoulli) head.} Support is the full power set, factorized into
independent per-type Bernoullis with $\rho_k(t)=\sigma(z_k)$:
\[
p^*(x\mid t)=\prod_{k\in x}\rho_k\prod_{k\notin x}(1-\rho_k),\qquad
\sum_{x\subseteq\{1,\dots,K\}}p^*(x\mid t)=\prod_{k=1}^{K}\bigl(\rho_k+(1-\rho_k)\bigr)=1.
\]
Both heads satisfy $\sum_x p^*(x\mid t)=1$, so $\sum_x\lambda_x^*(t)=\Lam(t)$ and the rate pin of
\S\ref{sec:pin} holds \emph{either way}. A single event is the singleton $x=\{k\}$ ($|x|=1$), so the
set-valued form subsumes single-item output as a special case. (For crypto LOBs simultaneity is
$<3\%$, so the categorical head is used.)

\section{Likelihood and training}
For a window $\Hist_T$ the log-likelihood is the standard MTPP objective
\[
\log\mathcal{L}=\sum_{i=1}^{N_T}\Bigl[\log\Lam(t_i)+\log p^*(x_i\mid t_i)\Bigr]-\int_0^T\Lam(u)\,du,
\]
where the compensator $\int_0^T\Lam$ has the closed form
$\mu_0 T+\sum_m \tfrac{a_m}{\delta_m}\sum_{t_i<T}\bigl(1-e^{-\delta_m(T-t_i)}\bigr)$.
The time and mark terms are optimized jointly; $\mu_0$ is recomputed from $n$ each step
(\S\ref{sec:pin}) rather than learned freely.

\section{Why this parameterization}
\begin{itemize}[leftmargin=1.4em,itemsep=2pt]
  \item \textbf{Calibration by construction.} The mean event rate is pinned exactly, so simulated
  rollouts cannot drift in activity --- the dominant failure mode of free-rate neural MTPPs.
  \item \textbf{Certifiable stability.} $n<1$ is checked in closed form every step; the process is
  provably subcritical, so long closed-loop simulation does not explode.
  \item \textbf{Exact Ogata thinning.} The linear ground rate is monotone non-increasing between
  events, giving an exact upper bound $\bar\lambda=\Lam(t_i^+)$ --- so LGM admits \emph{exact}
  Ogata thinning, unlike nonlinear neural compensators that need approximate grid inversion.
  \item \textbf{Clean ablation.} Running the same per-type mark network standalone with
  $\lambda_k=\softplus(z_k)$ (driving the rate itself) recovers the PCT-LSTM baseline; the
  LGM-vs-PCT-LSTM gap isolates the contribution of the linear rate-pinned ground rate.
\end{itemize}

\end{document}
